\chapter{Zusammenfassung und Ausblick} \label{zusammenfassung}
% [X] Übersicht der Teilkapitel
In diesem Kapitel wird eine Zusammenfassung der wichtigsten Ergebnisse der Arbeit erläutert. Anschließend wird ein Überblick über die nicht in dieser Arbeit betrachteten Funktionalitäten gegeben, sowie ein Ausblick auf mögliche Erweiterungen, die zur Verbesserung des Systems beitragen können.

\section{Fazit}
% [X] Hat zur POC gedient, und man kann erfolgreich Testen, in zwei Betriebsmodus durchzuführen usw...
Der Schwerpunkt dieser Arbeit lag auf der Erörterung und Entwicklung eines Konzeptes zur schnellen und einfachen Integration von Simulink-Modellen in bestehende BaSyx-Architekturen. Zu diesem Zweck wurden zwei Anwendungsfälle betrachtet: die Verbindung der BaSyx-Komponenten in Echtzeit sowie im Abbild-Modus für die asynchrone Bearbeitung der Daten der Komponenten.
% [X] Modulare Konzepierung, die leicht durch andere Modulen erweitbar ist

Das ganze System wurde modular entworfen, sodass künftige Erweiterungen der Funktionalität leichter implementierbar sind. Neue Blöcke und Anwendungsfälle können bzw. auf vorhandene Methoden zugreifen, um dadurch über eine getestete Abstraktion der für C++ vorhandenen Matlab-Logik zu verfügen.    

Außerdem wurde ein Testszenario zur Validierung des Erfüllens aller festgelegten Anforderungen konzipiert. Dieses Szenario zeigt weiterhin, wie eine übliche Anwendung des Maschinenbaus und der Verfahrenstechnik mittels Simulink und BaSyx bereitgestellt werden kann.

% [X] Daraus lässt sich die Schlussfolgerung ziehen, dass
Die Arbeit zeigt schließlich, dass durch die Anwendung der entwickelten Bibliothek \textit{BaSyx} und \textit{Simulink} unkompliziert miteinander verbunden werden können. Ihr Einsatz eignet sich besonders gut für das Testen von Komponenten in früheren Phasen des Produktlebenszyklus, um Prototypen aufzubauen und unterschiedliche Varianten eines Konzepts zu testen.


\section{Ausblick}
% [X] Anforderung bezüglich Geschwindikeit usw müssen noch berücksichtigt werden, Optimierungen
Da sich das Projekt sich prinzipiell mit dem strukturellen Testen der Anbindung von BaSyx und Simulink auseinandergesetzt hat, sollten Erweiterungen bezüglich der Sicherheit und Ausführungsdauer in späteren Iterationen berücksichtigt werden.

Der in dieser Arbeit entworfene Block reicht für eventuelle Anforderungen anderer Anwendungen nicht aus. Künftig könnten weitere Blöcke erstellt werden, die beispielsweise Text-Werte verarbeiten und den Kommunikationsfluss verbessern können.   

% [X] Einfügen von weiteren Blöcken, bzw. für die Verbindung mitm Server
Darüber hinaus wäre es sinnvoll, die Bibliothek in einer anspruchsvolleren Anwendung einzusetzen, um das Potential und die Vorteile einer Digitalisierungsstruktur mit Tools wie BaSyx und Simulink hervorzuheben sowie deren Grenzen zu verdeutlichen.

% [ ] @tagline Erwähnen dass die Lib erweitertet  und in eine AF getestet wird
